\documentclass[11pt,a4paper]{article}
\usepackage[margin=27mm]{geometry}
\usepackage{amsmath,amsthm,amssymb,fontspec,unicode-math,hyperref}
\usepackage{longtable,booktabs,array,calc,fvextra}
\setmainfont{STIXGeneral}
\setmathfont{STIX Two Math}
\setmonofont{Menlo}
\hypersetup{hidelinks}
\newtheorem{theorem}{Theorem}
\newtheorem{proposition}[theorem]{Proposition}
\newtheorem{lemma}[theorem]{Lemma}
\newtheorem{corollary}[theorem]{Corollary}
\theoremstyle{remark}\newtheorem{remark}[theorem]{Remark}
\newcommand{\gen}[1]{\langle #1\rangle}
\newcommand{\F}{\mathbb F}
\newcommand{\Z}{\mathbb Z}
\newcommand{\Aut}{\operatorname{Aut}}
\newcommand{\Sym}{\operatorname{Sym}}
\setlength{\emergencystretch}{2em}
\providecommand{\tightlist}{\setlength{\itemsep}{0pt}\setlength{\parskip}{0pt}}

\title{Recovering incidence from a complete neighbourhood}
\author{Human Compression Project}\date{}
\begin{document}\maketitle
\begin{abstract}
A transitive group action identifies objects with stabilizer cosets. The missing condition for reconstructing incidence is a converse: every incident pair must be in the diagonal orbit of a chosen base pair. We prove that this is equivalent to one complete neighbourhood calculation for each type pair, and give its finite-cardinality form.
\end{abstract}
\section{The complete reconstruction theorem}
Let \(G\) act on the right on nonempty sets \(X_i\), transitively on
each type. For distinct types, let \(R_{ij}\subseteq X_i\times X_j\) be
a \(G\)-invariant binary relation. Choose objects \(x_i\in X_i\), put
\(H_i=G_{x_i}\), and suppose \(x_i R_{ij}x_j\) for every pair of
distinct types.

\textbf{Theorem C.Incidence (equivariant reconstruction).} Suppose that, for
every \(i\ne j\), \[
 \{y\in X_j:x_iR_{ij}y\}=x_j^{H_i}.
 \tag{2.1}
\] Then the maps \[
 H_i g\longmapsto x_i^g
\] give an isomorphism from the right-coset incidence structure of the
subgroups \(H_i\) to the given typed incidence structure. If the
neighbour sets are finite, it is enough to check \[
 \#\{y\in X_j:x_iR_{ij}y\}
   =[H_i:H_i\cap H_j].
 \tag{2.2}
\]

\emph{Proof.} Equality of right cosets is equivalent to equality of the
associated images of \(x_i\), because \(H_i\) is exactly its stabilizer.
Transitivity gives surjectivity. Two cosets \(H_i g,H_j h\) intersect if
and only if there is \(a\in G\) with \(x_i^a=x_i^g\) and
\(x_j^a=x_j^h\). Thus coset incidence is exactly the diagonal orbit of
\((x_i,x_j)\).

The base pair lies in \(R_{ij}\), so invariance puts that entire
diagonal orbit inside \(R_{ij}\). Conversely, suppose \(x_i^gR_{ij}y\).
Applying \(g^{-1}\) gives \(x_iR_{ij}y^{g^{-1}}\). By (2.1),
\(y^{g^{-1}}=x_j^u\) for some \(u\in H_i\). Then
\((x_i^g,y)=(x_i^{ug},x_j^{ug})\), so the pair lies in the diagonal
orbit. This proves equality of the binary relations. All flags and
residues are defined by these relations, and therefore are preserved as
well.

For the finite criterion, invariance of \(R_{ij}\) and the base
incidence give \(x_j^{H_i}\subseteq R_{ij}(x_i)\). Orbit--stabilizer
gives \(|x_j^{H_i}|=[H_i:H_i\cap H_j]\). Equality of the two finite
cardinalities proves equality of the sets. \(\square\)

The hypothesis in (2.2) concerns the \emph{complete} neighbour set.
Counting a chosen orbit of neighbours would merely repeat
orbit--stabilizer and would prove nothing about possible additional
incidences. The omitted relation in {[}HHL12{]} gives a concrete example
of this distinction.

The three Lean declarations \texttt{reconstruction},
\texttt{localOrbit\_subset} and \texttt{reconstruction\_of\_card} in
\texttt{formal/EquivariantIncidence.lean} formalize the
diagonal-orbit argument and the finite-cardinality step, with left
actions. Passing to inverse coordinates gives the right-action
convention above. They do not formalize the named graph, the whole coset
construction, the J1 certificates, or either complete source theorem.


\section{Use with the current common theory}
For a marking of a finite group, the complete subgroup intersection and typed coset Helly tests provide the hypotheses of Topic C's UC.Geometry. The theorem above then transfers every flag, every residue and every type-preserving automorphism to a source structure whose objects were defined independently. In particular a count of the whole neighbour set is essential; the order of one chosen stabilizer orbit cannot exclude additional incidences. The Livingstone application contains a literal example of exactly this failure.

This is the standard orbit--stabilizer argument, reused from the complete reviewed phase179 proof and its independent phase184 review. It is not a historical-novelty claim. The three unchanged Lean declarations formalize the diagonal-orbit converse and finite-set equality step, not the full coset quotient or the finite group computations. The application retains the complete current Topic C umbrella explicitly; no archived manuscript is a hidden premise.
\begin{thebibliography}{9}
\bibitem{HHL12} M. I. Hartley, I. Hubard and D. Leemans, \emph{Regular 4-polytopes from the Livingstone graph of Janko's first group}, Journal of Algebraic Combinatorics 35 (2012), 193--214, DOI 10.1007/s10801-011-0300-x.
\end{thebibliography}
\end{document}
